Projected entangled-pair states extend matrix product states to
higher-dimensional lattices by contracting virtual pairs through local
maps~\cite{VerstraeteCirac2004PEPS}. Their local tensor structure makes them
natural candidates for describing ground states whose entanglement is
concentrated near region boundaries. In one dimension, quantitative
Schmidt-tail truncation estimates connect entanglement spectra with global
MPS approximation~\cite{VerstraeteCirac2006MPS}. Hastings established the
entropy area law for unique ground states of uniformly gapped spin chains
with bounded finite-range interactions and obtained polynomial-dimension
MPS approximations~\cite{Hastings2007}. The higher-dimensional question
requires a representation compatible across many intersecting cuts.

There are several distinct PEPS approximation questions. The global
formulation in \cite[Section~2.2, Questions~7--8]{CiracGarreRubioPerezGarcia2019}
asks whether a spectral gap ensures polynomial bond dimension at
inverse-polynomial Hilbert-space norm error. That formulation uses
translation-invariant finite-range Hamiltonians on a torus.
Theorem~\ref{thm:main} addresses the corresponding assertion for arbitrary
bounded on-site and nearest-neighbor interactions on an open square.
Local-observable formulations instead control expectations on fixed-size
regions; the stronger structural formulation in
\cite[Conjecture~2]{SchwarzBuerschaperEisert2017} also requires an injective
approximant with a gapped parent Hamiltonian. The theorem here concerns the
full vector and imposes no such structural requirements on the approximant.

A positive local approximation result is due to Huang
\cite[Theorem~3]{Huang2026Local}. For states satisfying a R\'enyi entropy
area law of fixed order $0<\alpha<1$ across every rectangle, he obtains
PEPS of bond dimension $e^{O(1/\delta)}$, independent of lattice size,
that approximate norm-one local-observable expectations to error $\delta$.
The observable support scale is held fixed in this bound. The R\'enyi
hypothesis is stronger than a von Neumann entropy area law, and the
conclusion controls local expectations rather than global vector error.

An entropy bound and a representation theorem are different statements.
For arbitrary chain-state families, bounded von Neumann block entropy
does not ensure efficient global MPS approximation
\cite{SchuchWolfVerstraeteCirac2008}. In higher dimensions, Ge and Eisert
constructed families satisfying strong area laws that have no efficient
classical tensor-network descriptions
\cite[Definition~7 and Corollary~9]{GeEisert2016}. Their description model
includes a bound on the classical complexity of tensor entries; it should
be distinguished from a bound on virtual dimension with arbitrary complex
entries. These examples concern general states, rather than the uniformly
gapped ground-state class in Theorem~\ref{thm:main}. They explain why one
must control compatible operations and their accumulated error, in addition
to the spectra of separate marginals.

Thermal tensor networks provide another approach. Hastings constructed
local spectral filters and tensor-network descriptions of Gibbs operators
\cite{Hastings2006Local}. Moln\'ar, Schuch, Verstraete, and Cirac refined the
thermal approximation and derived quasipolynomial-dimension ground-state
PEPS under a gap together with a low-energy density-of-states bound
\cite[Section~III.C]{MolnarSchuchVerstraeteCirac2015}. At inverse-polynomial
accuracy their bond bound is exponential in a squared logarithm of system
size. The low-energy counting condition is an additional hypothesis.
A different cooling and postselection construction approximates the unique
ground state of a gapped $d$-dimensional local Hamiltonian at the boundary
of a $(d+1)$-dimensional PEPS
\cite[Section~V]{SchuchWolfVerstraeteCirac2007Complexity}; transferring that
representation to the original lattice is a further problem.

Recent tensor-network results also distinguish the dimension and the
approximation task. Arad, Firanko, and Jain obtained two-dimensional
mutual-information bounds for locally gapped frustration-free systems and
polynomial-bond matrix product operator approximations to maximally mixed
ground states in one dimension
\cite[Corollaries~2.7 and~2.9]{AradFirankoJain2026}. Such entropy or
one-dimensional operator bounds do not provide the simultaneous
original-grid PEPS conclusion needed here.

Our analytic inputs are the entropy area law and the stronger collar
estimates of the companion \cite{RepoAreaLaw}. The present proof adapts its
nested-filter and two-family ideas, and then develops compatible changes
of ownership, compression of a distributed density network, and routing on
the prescribed grid. The representation theorem uses the full companion
interface recorded in Appendix~\ref{app:companion}, rather than the
cutwise entropy inequality alone. Its polynomial bond bound permits
arbitrary tensor entries and is an existence statement. The distinction
from efficient contraction is substantive: general PEPS contraction has
worst-case complexity obstructions
\cite{SchuchWolfVerstraeteCirac2007Complexity}, although that fact alone does
not classify the particular approximants constructed here.
